% Kapitel 3 --- Forschungsfrage und Hypothesen
% Budget: 1 Seite
% Quelle im Vault: 20 Recherche/Hypothesen 3.1

\section{Forschungsfrage und Hypothesen}
\label{sec:frage}

\subsection{Fragestellung}

\todo{Die Fragestellung aus dem Steckbrief aufnehmen und auf die eigene Operationalisierung
zuspitzen: Welche Sachverhalte werden in mehreren Quellen unterschiedlich dargestellt, und
wie zuverlässig lässt sich das automatisch finden?}

\subsection{Hypothesen}

\begin{description}[style=nextline]
  \item[H1 --- Abweichungen konzentrieren sich auf instabile Sachverhaltstypen]
  \todo{Formulierung, Begründung aus der Literatur, Operationalisierung, und was als
  Widerlegung gilt.}
  \item[H2 --- Sprachmodellgestützte Extraktion ist präzise, aber nicht vollständig]
  \todo{Formulierung, Bezug zur Validierungsliteratur \parencite{bavaresco2025llms},
  Operationalisierung über Precision und Recall je Verfahrensvariante.}
  \item[H3 --- Die meisten Abweichungen sind Zeitartefakte]
  \todo{Formulierung, Begründung über Content Drift und die Änderungsfassungen der
  Ordnung, Operationalisierung über die Abweichungstypologie.}
\end{description}

\subsection{Was als Widerlegung gilt}

\todo{Für jede Hypothese einen Satz: Welches Ergebnis würde sie kippen? Ohne diesen Absatz
ist die Arbeit nicht falsifizierbar, und das fällt in der Bewertung auf.}
